\section{Higher multiplicity and singular observation coordinates}
\label{sec:multiplicity}

\subsection{The finite heat-polynomial normal form}
Suppose a real event has exact spatial order \(m\geq2\):
\[
 \partial_x^jh(\tau_*,x_*)=0\quad(0\leq j<m),\qquad
 a_m=\partial_x^mh(\tau_*,x_*)\ne0.
\]
For fixed \(m\), define
\begin{equation}\label{eq:heat-poly}
 p_m(S,Y)=
 \sum_{j=0}^{\lfloor m/2\rfloor}
 \frac{m!}{j!(m-2j)!}S^jY^{m-2j}.
\end{equation}

\begin{theorem}[Finite-multiplicity heat geometry]\label{thm:higher}
Uniformly for \((S,Y)\) in compact sets,
\begin{equation}\label{eq:higher-rescale}
 h(\tau_*+\epsilon^2S,x_*+\epsilon Y)
 =\frac{a_m}{m!}\epsilon^m p_m(S,Y)+O(\epsilon^{m+1}).
\end{equation}
In a sufficiently small fixed complex neighborhood of \(x_*\), the
cluster consists of \(m\) simple real zeros immediately before the event.
Immediately after it, there are no real zeros if \(m\) is even, and
exactly one simple real zero if \(m\) is odd. The remaining zeros form
nonreal conjugate pairs; the total complex multiplicity stays \(m\).
\end{theorem}
\begin{proof}
The coefficient of \(s^jy^k\) in the local Taylor series is
\(\partial_x^{2j+k}h/(j!k!)\), by the heat equation. All terms of
parabolic weight \(2j+k<m\) vanish, and the weight-\(m\) terms give
\eqref{eq:heat-poly}. A finite Taylor remainder, as in
Section~\ref{sec:double}, proves \eqref{eq:higher-rescale}.

For completeness, the real-root property of \(q_m(Y)=p_m(-1,Y)\)
can be proved directly. The displayed coefficients give
\[
 q_0=1,\quad q_1=Y,\quad
 q_{m+1}=Yq_m-2m q_{m-1}.
\]
These monic polynomials have simple real interlacing zeros by induction.
At successive zeros of \(q_m\), the values of \(q_{m+1}=-2m q_{m-1}\)
alternate in sign, giving \(m-1\) intervening roots. At the largest
zero of \(q_m\), \(q_{m-1}>0\) and \(q_{m+1}<0\), whereas
\(q_{m+1}(Y)>0\) for large positive \(Y\), giving a further root.
At the smallest zero, the sign of \(q_{m-1}\) is \((-1)^{m-1}\);
the sign of \(q_{m+1}\) there is \((-1)^m\), opposite its sign at
large negative \(Y\). This supplies the last root. The degree is
\(m+1\), so all are simple and the interlacing is strict.
The induction begins with \(q_1=Y\) and \(q_2=Y^2-2\).
Parity shows that precisely the odd polynomials have a zero at zero.

Set \(s=-\epsilon^2\) before the event. Dividing the field by
\(\epsilon^m\) gives an analytic family in \((\epsilon,Y)\), including
\(\epsilon=0\), with limiting simple roots those of \(q_m\).
The implicit-function theorem continues every root as a real simple
root for small real positive \(\epsilon\).
After the event, \(s=\epsilon^2\), and
\[
 p_m(1,\ii Y)=\ii^m p_m(-1,Y).
\]
The nonzero limiting roots are therefore purely imaginary and stay
nonreal under a sufficiently small perturbation. If \(m\) is odd,
the limiting root at zero is simple and its unique continued branch
is real by conjugation symmetry. For even \(m\) there is no such root.
Finally, choose a fixed disk containing only the order-\(m\) zero at
the event. Its nonzero boundary and analytic convergence preserve the
total zero count \(m\). The branches just constructed exhaust it.
\end{proof}

This recovers the local forward-parabolic multiplicity pattern in
\cite{Angenent1991,HeatMultiplicity}. It asserts neither that all
canonical events have multiplicity two nor that one can choose the neighborhoods
or remainder constants uniformly in \(m\) or in height.

\subsection{What the generator detects at a node}
The complete operator intertwining gives, for each fixed \(j\),
\begin{equation}\label{eq:generator-jets}
 R_\chi K^j\Psi_\tau=(-1)^j\partial_x^{2j}h(\tau,\cdot).
\end{equation}
If \(m=2r\), the first nonzero pointwise value in this even-derivative
sequence is at \(j=r\). If \(m=2r+1\), the values for \(j\leq r\)
vanish, but the first spatial derivative of the \(j=r\) readout is
nonzero. These are justified pointwise statements for the prepared
smooth orbit. They do not assert bounded point evaluation on the
ambient \(L^2\) space. In particular, \(K\) and its powers are finite
on the source state even when lower-order readouts vanish.

\subsection{Logarithmic derivatives and their poles}
On an open region where \(h\ne0\), put \(w=h_x/h\).
Since \(h_\tau=h_{xx}\),
\[
 \partial_\tau\log h=\frac{h_{xx}}h=w_x+w^2,
\]
locally using any analytic logarithm. Differentiating in \(x\) gives
\begin{equation}\label{eq:burgers}
 w_\tau=w_{xx}+2ww_x.
\end{equation}
At a collision the local factorization gives
\[
 w(\tau_*,x_*+y)=\frac2y+\frac{b}{3a}+O(y),\qquad
 \frac{h_{xx}}h(\tau_*,x_*+y)=\frac2{y^2}+O(y^{-1}).
\]
These poles are created by division by the vanishing amplitude.
They are compatible with the regular field, regular generator
derivatives and regular energy identities already proved. Thus a
singularity in a logarithmic or root coordinate must be distinguished
from a singularity of the state.

All two-sided local statements concern the analytic source family
near an interior event. At the endpoint \(\tau=0\), a physical
semigroup orbit provides only the forward side. The entire source
formula still supplies a mathematical local extension, but it does
not make reverse evolution bounded on arbitrary ambient states.
